% GENERATED FILE. Edit canonical lessons, then run npm run generate:books.
% canonical-source-hash: fnv1a64:c9ea230db520c938
% canonical-lessons: GE-C326-vorheizen, GE-C326-einfetten, GE-C326-backblech, GE-C326-auflaufform, GE-C326-braeunen

\chapter{Ab in den Ofen}
\label{ch:ge-ab-in-den-ofen}
\hlchaptermodality{326}

\begin{chapteropening}
\textbf{By the end of this chapter:} \emph{I can describe getting ready to bake: to preheat, to grease, a baking tray, a baking dish and to brown.}

Complete the last lesson of chapter 326: I can describe getting ready to bake: to preheat, to grease, a baking tray, a baking dish and to brown.
\end{chapteropening}

\section[vorheizen]{vorheizen — to preheat}

\label{lesson:GE-C326-vorheizen}



\begin{quote}
Before the new one: say the German for to roll out, then the German for the rolling pin.
\end{quote}

\subsection*{You'll want to know: vorheizen}
\textbf{vorheizen} — \textquotedblleft{}to preheat\textquotedblright{}.

\begin{quote}
Er heizte den Ofen auf hundertachtzig Grad vor.
\end{quote}

\begin{quote}
\emph{He preheated the oven to a hundred and eighty degrees.}
\end{quote}

Separable: \emph{er heizt vor}, \emph{er hat vorgeheizt}.

\begin{cousinweb}
\emph{vorheizen} is built on \emph{heizen} \textquotedblleft{}to heat\textquotedblright{}, from Old High German \emph{heizen}, Proto-Germanic \emph{haitijaną} \textquotedblleft{}to make hot\textquotedblright{} (EtymDB), with \emph{vor} \textquotedblleft{}before\textquotedblright{}.
\end{cousinweb}

\begin{grammarlens}[title={What you've built}]
One more word for telling, explaining and describing at B1.
\end{grammarlens}

\subsection*{Your turn}
\begin{itemize}
\raggedright
  \item \emph{Say it:} \emph{vorheizen}
  \item \emph{Say it:} \emph{vorheizen}, once more
  \item \emph{Say it:} a sentence of your own with \emph{vorheizen}
  \item \emph{Recall:} say the German for to roll out, then the German for the rolling pin, then say \emph{vorheizen} again
\end{itemize}

\begin{tcolorbox}[breakable,colback=teal!4,colframe=teal!35!black,title={Before you move on}]
What does vorheizen mean? (\textquotedblleft{}To preheat\textquotedblright{}.) Say it once more.
\end{tcolorbox}
\section[einfetten]{einfetten — to grease}

\label{lesson:GE-C326-einfetten}



\begin{quote}
Before the new one: say the German for the rolling pin, then the German for to preheat.
\end{quote}

\subsection*{You'll want to know: einfetten}
\textbf{einfetten} — \textquotedblleft{}to grease\textquotedblright{}.

\begin{quote}
Sie fettete die Form mit Butter ein.
\end{quote}

\begin{quote}
\emph{She greased the tin with butter.}
\end{quote}

Separable. Built on \emph{das Fett}, fat.

\begin{cousinweb}
\emph{einfetten} is built on \emph{Fett} \textquotedblleft{}fat\textquotedblright{}, which came from Middle Low German \emph{vet} (EtymDB), a cousin of English \emph{fat}.
\end{cousinweb}

\begin{grammarlens}[title={What you've built}]
One more word for telling, explaining and describing at B1.
\end{grammarlens}

\subsection*{Your turn}
\begin{itemize}
\raggedright
  \item \emph{Say it:} \emph{einfetten}
  \item \emph{Say it:} \emph{einfetten}, once more
  \item \emph{Say it:} a sentence of your own with \emph{einfetten}
  \item \emph{Recall:} say the German for the rolling pin, then the German for to preheat, then say \emph{einfetten} again
\end{itemize}

\begin{tcolorbox}[breakable,colback=teal!4,colframe=teal!35!black,title={Before you move on}]
What does einfetten mean? (\textquotedblleft{}To grease\textquotedblright{}.) Say it once more.
\end{tcolorbox}
\section[das Backblech]{das Backblech — the baking tray}

\label{lesson:GE-C326-backblech}



\begin{quote}
Before the new one: say the German for to preheat, then the German for to grease.
\end{quote}

\subsection*{You'll want to know: das Backblech}
\textbf{das Backblech} — \textquotedblleft{}the baking tray\textquotedblright{}.

\begin{quote}
Die Kekse lagen auf dem Backblech.
\end{quote}

\begin{quote}
\emph{The biscuits lay on the baking tray.}
\end{quote}

Plural: \emph{die Backbleche}.

\begin{cousinweb}
\emph{Backblech} joins \emph{backen} \textquotedblleft{}to bake\textquotedblright{}, a cousin of English \emph{bake} (EtymDB), and \emph{Blech} \textquotedblleft{}sheet metal\textquotedblright{}, an inherited Germanic word.
\end{cousinweb}

\begin{grammarlens}[title={What you've built}]
One more word for telling, explaining and describing at B1.
\end{grammarlens}

\subsection*{Your turn}
\begin{itemize}
\raggedright
  \item \emph{Say it:} \emph{das Backblech}
  \item \emph{Say it:} \emph{das Backblech}, once more
  \item \emph{Say it:} a sentence of your own with \emph{das Backblech}
  \item \emph{Recall:} say the German for to preheat, then the German for to grease, then say \emph{das Backblech} again
\end{itemize}

\begin{tcolorbox}[breakable,colback=teal!4,colframe=teal!35!black,title={Before you move on}]
What does das Backblech mean? (\textquotedblleft{}The baking tray\textquotedblright{}.) Say it once more.
\end{tcolorbox}
\section[die Auflaufform]{die Auflaufform — the baking dish}

\label{lesson:GE-C326-auflaufform}



\begin{quote}
Before the new one: say the German for to grease, then the German for the baking tray.
\end{quote}

\subsection*{You'll want to know: die Auflaufform}
\textbf{die Auflaufform} — \textquotedblleft{}the baking dish\textquotedblright{}.

\begin{quote}
Er gab die Nudeln in eine Auflaufform.
\end{quote}

\begin{quote}
\emph{He put the pasta into a baking dish.}
\end{quote}

Plural: \emph{die Auflaufformen}. \emph{Der Auflauf} — a bake, a gratin.

\begin{cousinweb}
\emph{Auflaufform} joins \emph{Auflauf} \textquotedblleft{}a dish that rises\textquotedblright{}, from \emph{laufen} \textquotedblleft{}to run\textquotedblright{} (EtymDB), and \emph{Form} \textquotedblleft{}shape, tin\textquotedblright{}, from Latin \emph{forma} \textquotedblleft{}shape\textquotedblright{} (Lewis \& Short).
\end{cousinweb}

\begin{grammarlens}[title={What you've built}]
One more word for telling, explaining and describing at B1.
\end{grammarlens}

\subsection*{Your turn}
\begin{itemize}
\raggedright
  \item \emph{Say it:} \emph{die Auflaufform}
  \item \emph{Say it:} \emph{die Auflaufform}, once more
  \item \emph{Say it:} a sentence of your own with \emph{die Auflaufform}
  \item \emph{Recall:} say the German for to grease, then the German for the baking tray, then say \emph{die Auflaufform} again
\end{itemize}

\begin{tcolorbox}[breakable,colback=teal!4,colframe=teal!35!black,title={Before you move on}]
What does die Auflaufform mean? (\textquotedblleft{}The baking dish\textquotedblright{}.) Say it once more.
\end{tcolorbox}
\section[bräunen]{bräunen — to brown}

\label{lesson:GE-C326-braeunen}



\begin{quote}
Before the new one: say the German for the baking tray, then the German for the baking dish.
\end{quote}

\subsection*{You'll want to know: bräunen}
\textbf{bräunen} — \textquotedblleft{}to brown\textquotedblright{}.

\begin{quote}
Der Käse bräunte im Ofen.
\end{quote}

\begin{quote}
\emph{The cheese browned in the oven.}
\end{quote}

Also of skin in the sun: \emph{sich in der Sonne bräunen}.

\begin{cousinweb}
\emph{bräunen} is built on \emph{braun} \textquotedblleft{}brown\textquotedblright{}, Old High German \emph{brūn}, Proto-Germanic \emph{brūnaz} (EtymDB).
\end{cousinweb}

\begin{grammarlens}[title={What you've built}]
One more word for telling, explaining and describing at B1.
\end{grammarlens}

\subsection*{Your turn}
\begin{itemize}
\raggedright
  \item \emph{Say it:} \emph{bräunen}
  \item \emph{Say it:} \emph{bräunen}, once more
  \item \emph{Say it:} a sentence of your own with \emph{bräunen}
  \item \emph{Recall:} say the German for the baking tray, then the German for the baking dish, then say \emph{bräunen} again
\end{itemize}

\begin{tcolorbox}[breakable,colback=teal!4,colframe=teal!35!black,title={Before you move on}]
What does bräunen mean? (\textquotedblleft{}To brown\textquotedblright{}.) Say it once more.
\end{tcolorbox}
